\chapter{Netteté}
\section{Les trois contrôles}
Le \emph{Gain de netteté} va de 0 à 200 et règle la force du renforcement.
Le \emph{Rayon de netteté} va de 0,3 à 3,0 et règle son échelle ; sa référence
est 1,0. Le \emph{Masquage de netteté} va de 0 à 100 et module la sélectivité
du renforcement. Gain et masquage ont une référence à zéro.

\section{Procédure}
Évaluer d’abord le bruit, puis augmenter légèrement le gain de netteté.
Faire varier le rayon avec prudence et observer les contours contrastés.
Rechercher un gain de lisibilité sans liseré visible ni accentuation des
zones uniformes. Le renforcement ne récupère pas une mise au point réellement absente.

\paragraph{Exemple.}
Sur un portrait, contrôler le bord de l’œil et la texture de la peau, pas
uniquement les cheveux. Basculer entre \touche{Ctrl + 1} et \touche{Ctrl + 0}.
Maintenir \touche{Espace} pour comparer à la référence ; contrôler ensuite
un export à la taille finale.
\gestesreglages
